
\section{Introduction}
\label{sec:intro}

The removal of unwanted information is a surprisingly common task. Transform-invariant features in computer vision, ``fair'' encodings from the algorithmic fairness community, and two-stage regressions often used in scientific studies are all cases of the same general concept:
we wish to remove the effect of some outside variable $c$ on our data $x$ while still relevant to our original task. In the context of representation learning, we wish to map $x$ into an encoding $z$ that is uninformative of $c$, yet also optimal for our task loss $\mathcal{L}(z,\dots)$. %In otherwords, we would like a representation $z$ of data $x$ that is invariant to an outside factor $c$. %These $z$ are invariant to $c$
%we wish to perform tasks using data $x$ without biases toward outside variable $c$, of which $x$ is informative. The outside variable $c$ might be e.g. rotations of an image, information on the protected classes of users (gender, race, etc.), or equipment data or other meta-experimental information.

%we wish to map data $x$ into an encoding $z$ that contains little or no information about another variable $c$, while keeping $z$ relevant to our original task.

%we wish to remove information about an outside variable $c$ fr
%we wish to find a representation $z$ of our data $x$ that is at once both useful for our original task yet also minimally informative about an outside $c$. In our 
  %We further desire $z$ to be descriptive either of a label $y$ or the data $x$ itself (the supervised and unsupervised cases respectively).

%Put into representation learning terms, we wish to map data $x$ into an encoding $z$ that is uninformative of $c$ yet still relevant to our original task.
%These qualitative objectives suggest a general quantitative criterion for learning such $z$,
%that we should use the mutual information $I(z,c)$ to penalize representations informative about $c$.
%Following the intuition of these qualitative objectives, various penalties and constraints have been proposed to encourage or enforce a $z\perp c$ condition.
These objectives are often operationalized as an independence constraint $z\perp c$.
Encodings satisfying this condition are invariant under changes in $c$, thus called ``invariant representations''.
In practice these constraints are often relaxed to other measures; 
%seek $z$ that is minimally informative about $c$. In other words, we should use the mutual information $I(z,c)$ to penalize representations $z$ which contain more information about $c$.
%In combination with the task objective (practically, another loss function), we should seek $z$ that is minimally informative about $c$. In other words, we should use the mutual information, $I(z,c)$, to penalize representations $z$ which contain more information about $c$.
in recent works an adversary's ability to predict $z$ from $c$ has been used as such a proxy \cite{louizos2015variational,xie2017controllable}, transplanting adversarial losses from generative literature to encoder/decoder settings.

In the present work we instead relax $z\perp c$ to a penalty on the mutual information $I(z,c)$. We provide an analysis of this loss, showing that:
\begin{enumerate}
\item $I(z,c)$ admits a useful variational upper bound. This is in contrast to the usual lower bound e.g. bounds on $I(z,y)$ for some labels $y$.
\item When placed alongside the Variational Auto Encoder (VAE) and the Variational Information Bottleneck (VIB) frameworks, the upper bound on $I(z,c)$ produces a computationally tractable form for learning $c$-agnostic encodings and predictors.
\item The adversarial approach can be also derived as a procedure to minimize $I(z,c)$, but does not provide an upper bound.
\end{enumerate}

Our proposed methods have the practical advantage of only requiring $c$ at training time, and not at test time. They are thus viable for production settings where accessing $c$ is expensive (requiring human labeling), impossible (requiring underlying transformations), or legally inadvisable (sharing protected data). On the other hand, our method also produces a conditional decoder taking both $c$ and $z$ as inputs. While for some purposes this might be discarded at test time, it also can be manipulated to function similarly to Fader Networks~\cite{lample2017fader}, in that we can generate realistic looking transformations of an input image where some class label has been altered. We empirically test our proposed $c$-agnostic VAE and VIB methods on two standard ``fair prediction'' tasks, as well as an unsupervised learning task demonstrating Fader-like capabilities.

\subsection{Related Work}

The removal of covariate factors from scientific data has a long history. Observational studies in the sciences often cannot control for every factor influencing subjects, thus a large amount of literature has been generated on the topic of removing such factors after data collection. Simple statistical techniques usually involve modifying analyses with corresponding covariate effects \cite{raudenbush1994random} or sometimes multi-level regressions \cite{freckleton2002misuse}. Modern methods have included more nuanced feature generation and more complex models, but follow along the same vein \cite{feis2015ica,fortin2017harmonization}
``Regressing out'' or ``controlling for'' unwanted factors can be effective, but place strong constraints on later analyses (namely, the observation of the same unwanted covariates).%Generalization to unseen covariate effects is usually not considered.

A similar concept also has deep roots in computer vision, where transform-invariant features and/or methods have been sought after for some time. Often these methods were designed for specific cases, e.g. scale-invariant or rotation invariant features. Early examples include Steerable Filters \cite{freeman1991design,greenspan1994overcomplete}, and later SIFT \cite{lowe1999object}. For many image transformations, data augmentation has become a standard practical tool to encourage invariance.  

Recent work has provided a group-theoretic \cite{cohen2016group} analysis of the removal of covariate information (either by design or by augmentation), in which equivalences are drawn between finding invariant features and finding the quotient space of the domain over a covariate group action. More generally, an empirical solution was proposed by Lample et al. \cite{lample2017fader}, who removed specific visual features from a latent representation through adversarial training. 

More recently the algorithmic fairness community has investigated fair methods \cite{kamiran2009classifying} and fair representations \cite{zemel2013learning}. Derived in part from the desire to avoid discriminating against protected classes of individuals (and in part to avoid breaking laws and/or to avoid being the subject of a civil suit), the objective of these methods has been to preserve task accuracy (usually classification or regression) while removing bias against the protected class of individuals. 

Particularly relevant to our work are the recent methods proposed by Louizos et al. \cite{louizos2015variational} and Xie et al. \cite{xie2017controllable}, which have a similar problem setup. Both methods generate representations that make it difficult for an adversary to recover the protected class but are still useful for a classification task. Louizos et al. propose the ``Variational Fair Auto-Encoder'' (VFAE), which, as its name suggests, modifies the VAE of Kingma and Welling \cite{kingma2013auto} to produce fair\footnote{The definition of ``fair'' in an algorithmic setting is of some debate. A fair encoding in this paper is uninformative of protected classes. We offer no opinion on whether this is truly ``fair'' in a general or legal sense, taking the word at face value as used by Louizos et al.} encodings, as well as a supervised case providing fair classifications. Xie et al. combine this concept with adversarial training, adding (inverted) gradient information to produce fair representations and classifications. This adversarial solution coincides exactly with the conceptual framework used in a computer vision application for constructing Fader Networks \cite{lample2017fader}.

Compressive encoding in a learning context is also well studied. In particular, the Information Bottleneck \cite{tishby2000information} and its modern successor Variational Information Bottleneck (VIB) \cite{alemi2016deep,achille2018information} both provide compressive encodings, aiming for ``relevance'' with respect to a target variable (usually a label). An unsupervised method, CorEx, also has a similar extension (Anchored Corex \cite{gallagher2016anchored}), in which latent factors can be driven toward specific targets. Our work could be thought of as adding ``negative anchors'' or aiming for ``irrelevance'' with respect to protected classes.

Models including ``nuisance'' factors were also considered by Soatto and Chiuso \cite{soatto2014visual}, in which the authors propose definitions of both nuisances and invariant representations. The authors follow the group theoretic concept of nuisances. Achille and Soatto \cite{achille2017emergence} directly utilize these results, providing the same criterion and relaxation for invariance we will use here, minimal mutual information between the representation $z$ and the covariate $c$. While nuisances form only a small subsection of the paper, the authors propose and test a sampling based approach to learning invariant representation. Their method is predicated on the ability to sample from the nuisance distribution (e.g. adding occlusions in images). We optimize a similar objective, but avoid such practical constraints.

Contemporaneous to this work, another paper by Song et al. \cite{song2018learning} proposes a similar information theory based bound. Instead of making the variational approximation of reconstruction using $p(x|z,c)$, the authors give the following bound $I(z,c) \leq \mathbb{E}[~KL[q(z|x,c) \| p(z) ]$ for a specified $p(z)$. This leads to a similar desiderata, which they optimizing using an (adversarial) iterative minimax method.


